\documentclass[11pt]{article}
\usepackage[margin=1in]{geometry}
\usepackage{xcolor}
\usepackage[breakable,listings]{tcolorbox}
\tcbuselibrary{listings}
\usepackage{hyperref}
\usepackage{enumitem}
\usepackage{listings}
\usepackage{amsmath}
\usepackage{booktabs}
\usepackage{array}

% Define colors
\definecolor{osaorange}{RGB}{220,115,38}
\definecolor{aiblue}{RGB}{30,144,255}

% Configure hyperref
\hypersetup{
    colorlinks=true,
    linkcolor=blue,
    urlcolor=blue
}

% Code environment
\newtcblisting{rcode}{
  colback=white,
  colframe=gray!50,
  boxrule=0.5pt,
  arc=0pt,
  left=3pt,
  right=3pt,
  top=3pt,
  bottom=3pt,
  width=\textwidth,
  breakable,
  listing only,
  listing options={
    basicstyle=\small\ttfamily,
    breaklines=true,
    breakatwhitespace=true,
    columns=flexible,
    keepspaces=true,
    showstringspaces=false
  }
}

\title{}
\author{}
\date{}

\begin{document}

% Header box
\begin{tcolorbox}[colback=osaorange,colframe=black,boxrule=2pt,arc=0pt,left=3pt,right=3pt,top=6pt,bottom=6pt,boxsep=2pt]
\centering
\textcolor{white}{\textbf{\Large H524 Introduction to Biostatistics}}\\
\textcolor{white}{\textbf{\Large Week 2: Introduction to Probability}}\\
\textcolor{white}{\textbf{\Large Homework Assignment}}\\
\textcolor{white}{\textbf{\Large Fall 2025}}
\end{tcolorbox}

\vspace{8pt}

\noindent\textbf{Due Date:} End of Week 3\\
\textbf{Total Points:} 100 (+ 5 bonus)\\
\textbf{Format:} Submit through Canvas

\section*{Instructions}

\begin{itemize}
\item Show all work for full credit
\item Round probabilities to 4 decimal places unless otherwise specified
\item Round percentages to 1 decimal place
\item Include R code where specified
\item You may use R to check your work
\end{itemize}

\section{Part 1: Basic R and Probability Fundamentals (25 points)}

\subsection*{Question 1: Basic R Operations (5 points)}

Write R code to:

\begin{enumerate}[label=\alph*)]
\item Create a vector with the following systolic blood pressures: 118, 132, 145, 128, 155, 122, 138 (1 point)
\item Calculate the mean and standard deviation (2 points)
\item Determine how many patients have high blood pressure ($\geq$ 140 mmHg) (2 points)
\end{enumerate}

\vspace{1.5cm}

\subsection*{Question 2: Simple Probability (5 points)}

In a hospital emergency department, records show:
\begin{itemize}
\item 40\% of patients come for injuries
\item 35\% come for illnesses
\item 25\% come for other reasons
\end{itemize}

\begin{enumerate}[label=\alph*)]
\item What is the probability the next patient does NOT come for an injury? (2 points)
\item If 200 patients visit the ED this week, how many would you expect to come for illnesses? (3 points)
\end{enumerate}

\vspace{1.5cm}

\subsection*{Question 3: Sample Space and Events (5 points)}

A clinical trial randomly assigns patients to three treatment groups: Drug A, Drug B, or Placebo.
\begin{itemize}
\item 50 patients receive Drug A
\item 75 patients receive Drug B
\item 75 patients receive Placebo
\end{itemize}

\begin{enumerate}[label=\alph*)]
\item What is the probability a randomly selected patient received Drug A? (2 points)
\item What is the probability a randomly selected patient did NOT receive placebo? (3 points)
\end{enumerate}

\vspace{1.5cm}

\subsection*{Question 4: Complement Rule (5 points)}

A screening test for diabetes has a 95\% specificity (true negative rate).

\begin{enumerate}[label=\alph*)]
\item What is the probability of a false positive? Show your work using the complement rule. (3 points)
\item Explain in plain language what this false positive probability means. (2 points)
\end{enumerate}

\vspace{1.5cm}

\subsection*{Question 5: Addition Rule (5 points)}

In a study of 300 adults:
\begin{itemize}
\item 120 have high blood pressure
\item 90 have high cholesterol
\item 30 have both conditions
\end{itemize}

\begin{enumerate}[label=\alph*)]
\item Calculate $\Pr(\text{high BP OR high cholesterol})$ using the addition rule. Show your work. (4 points)
\item Are these events mutually exclusive? Explain. (1 point)
\end{enumerate}

\vspace{1.5cm}

\section{Part 2: Conditional Probability (20 points)}

\subsection*{Question 6: Basic Conditional Probability (6 points)}

In a group of 500 hospital patients:
\begin{itemize}
\item 200 are smokers
\item 150 have respiratory complications
\item 100 are smokers with respiratory complications
\end{itemize}

\begin{enumerate}[label=\alph*)]
\item Calculate $\Pr(\text{respiratory complications} \mid \text{smoker})$. Show your work. (3 points)
\item Calculate $\Pr(\text{smoker} \mid \text{respiratory complications})$. Show your work. (3 points)
\end{enumerate}

\vspace{1.5cm}

\subsection*{Question 7: Law of Total Probability (7 points)}

A clinic's patient population:
\begin{itemize}
\item 60\% are adults, 40\% are children
\item Among adults, 20\% have allergies
\item Among children, 35\% have allergies
\end{itemize}

\begin{enumerate}[label=\alph*)]
\item What is the overall probability that a randomly selected patient has allergies? Show your work using the law of total probability. (5 points)
\item If a patient has allergies, what is the probability they are a child? (2 points)
\end{enumerate}

\vspace{1.5cm}

\subsection*{Question 8: Independence (7 points)}

Two screening tests are given:
\begin{itemize}
\item $\Pr(\text{Test A positive}) = 0.10$
\item $\Pr(\text{Test B positive}) = 0.15$
\item $\Pr(\text{both positive}) = 0.015$
\end{itemize}

\begin{enumerate}[label=\alph*)]
\item Are Test A and Test B independent? Show your work. (5 points)
\item Explain what independence would mean in this clinical context. (2 points)
\end{enumerate}

\vspace{1.5cm}

\section{Part 3: Diagnostic Testing and 2$\times$2 Tables (30 points)}

\subsection*{Question 9: Sensitivity and Specificity Definitions (6 points)}

A rapid flu test is being evaluated.

\begin{enumerate}[label=\alph*)]
\item Define sensitivity in terms of conditional probability using proper notation [$\Pr(\cdot \mid \cdot)$]. (2 points)
\item Define specificity in terms of conditional probability using proper notation. (2 points)
\item Explain in plain language what a 90\% sensitivity means for patients who actually have the flu. (2 points)
\end{enumerate}

\vspace{1.5cm}

\subsection*{Question 10: Creating a 2$\times$2 Table (10 points)}

A COVID-19 rapid antigen test has:
\begin{itemize}
\item Sensitivity = 85\%
\item Specificity = 95\%
\end{itemize}

In a community where 2\% of people currently have COVID-19, 10,000 people are tested.

\begin{enumerate}[label=\alph*)]
\item Create a complete 2$\times$2 table showing:
   \begin{itemize}
   \item Disease status (rows): Has COVID / No COVID
   \item Test result (columns): Positive / Negative
   \item Fill in all four cells with actual counts (6 points)
   \end{itemize}
\item Add row and column totals to your table. (2 points)
\item Verify your sensitivity and specificity calculations match the given values. (2 points)
\end{enumerate}

\vspace{3cm}

\subsection*{Question 11: PPV and NPV Calculation (8 points)}

Using the 2$\times$2 table from Question 10:

\begin{enumerate}[label=\alph*)]
\item Calculate the Positive Predictive Value (PPV). Show your work. (3 points)
\item Calculate the Negative Predictive Value (NPV). Show your work. (3 points)
\item Explain in one sentence what the PPV means for someone who tests positive. (2 points)
\end{enumerate}

\vspace{1.5cm}

\subsection*{Question 12: The Counterintuitive Case (6 points)}

Consider an HIV screening test with excellent performance:
\begin{itemize}
\item Sensitivity = 99.5\%
\item Specificity = 98.5\%
\item Population prevalence = 0.1\% (1 in 1,000)
\end{itemize}

Imagine screening 100,000 people.

\begin{enumerate}[label=\alph*)]
\item Without doing calculations, would you expect the PPV to be high or low? Why? (2 points)
\item Create the 2$\times$2 table with actual counts for 100,000 people screened. (3 points)
\item Calculate the PPV and explain why it might surprise people despite the excellent sensitivity and specificity. (1 point)
\end{enumerate}

\vspace{3cm}

\section{Part 4: Risk Measures (15 points)}

\subsection*{Question 13: Relative Risk (8 points)}

A 5-year cohort study of 1,000 people examined coffee consumption and heart disease:
\begin{itemize}
\item 400 heavy coffee drinkers: 60 developed heart disease
\item 600 light/non coffee drinkers: 48 developed heart disease
\end{itemize}

\begin{enumerate}[label=\alph*)]
\item Calculate the risk of heart disease in heavy coffee drinkers. (2 points)
\item Calculate the risk of heart disease in light/non coffee drinkers. (2 points)
\item Calculate the Relative Risk (RR). (2 points)
\item Interpret the RR in one complete sentence. (2 points)
\end{enumerate}

\vspace{1.5cm}

\subsection*{Question 14: Understanding Risk Ratios (7 points)}

For each Relative Risk value below, explain what it means:

\begin{enumerate}[label=\alph*)]
\item RR = 1.0 (2 points)
\item RR = 2.5 (2 points)
\item RR = 0.6 (2 points)
\item When RR $>$ OR, is the disease common or rare? (1 point)
\end{enumerate}

\vspace{1.5cm}

\section{Part 5: Probability Distributions in R (10 points)}

\subsection*{Question 15: Binomial Distribution (5 points)}

A new medication has a 75\% success rate. A hospital treats 20 patients with this medication.

\begin{enumerate}[label=\alph*)]
\item What is the probability that exactly 15 patients respond successfully? Write the R code and provide the answer. (2 points)
\item What is the probability that at least 16 patients respond successfully? Write the R code and provide the answer. (3 points)
\end{enumerate}

\vspace{1.5cm}

\subsection*{Question 16: Normal Distribution (5 points)}

Adult cholesterol levels are normally distributed with mean = 200 mg/dL and SD = 40 mg/dL.

\begin{enumerate}[label=\alph*)]
\item What proportion of adults have cholesterol levels above 240 mg/dL (high risk)? Write the R code and provide the answer. (2 points)
\item What cholesterol level represents the 90th percentile? Write the R code and provide the answer. (3 points)
\end{enumerate}

\vspace{1.5cm}

\section*{Bonus: AI-Enhanced Learning (5 bonus points)}

\subsection*{Question 17: AI Verification (5 bonus points)}

Choose ONE problem from Part 3 (Diagnostic Testing).

\begin{enumerate}[label=\alph*)]
\item Ask an AI tool (like Claude, ChatGPT) to solve it. Include your exact prompt. (1 point)
\item Verify the AI's answer by working through it yourself. Did the AI get it right? (2 points)
\item Identify one strength and one potential limitation of using AI for this type of problem. (2 points)
\end{enumerate}

\vspace{2cm}

\section*{Submission Requirements}

\subsection*{Required Files:}
\begin{enumerate}
\item \textbf{PDF Document} with all written work, tables, and calculations
\item \textbf{R Script File (.R)} with all R code properly commented
\item \textbf{AI Usage Documentation} (if applicable):
   \begin{itemize}
   \item Prompts used
   \item How you verified outputs
   \item Your own understanding of the solutions
   \end{itemize}
\end{enumerate}

\subsection*{Formatting Guidelines:}
\begin{itemize}
\item Clearly label each question number
\item Show all calculation steps
\item For 2$\times$2 tables, draw clear grids with labels
\item Round consistently per instructions
\end{itemize}

\section*{Grading Rubric}

\begin{table}[h]
\centering
\begin{tabular}{lcp{8cm}}
\toprule
\textbf{Category} & \textbf{Points} & \textbf{Criteria} \\
\midrule
Calculations & 40 & Correct formulas, accurate arithmetic \\
Methods & 30 & Appropriate approach, proper notation \\
Interpretation & 20 & Clear explanations, correct context \\
R Code & 10 & Working code, proper comments \\
Bonus & +5 & Thoughtful AI integration \\
\bottomrule
\end{tabular}
\end{table}

\section*{Academic Integrity}

\begin{tcolorbox}[colback=aiblue!10,colframe=aiblue,title=AI Usage Guidelines]

\textbf{You are encouraged to:}
\begin{itemize}
\item[+] Use AI tools to check your work
\item[+] Ask AI to explain concepts
\item[+] Verify calculations with AI
\end{itemize}

\textbf{You must NOT:}
\begin{itemize}
\item[-] Submit AI-generated work without understanding it
\item[-] Copy solutions without showing your reasoning
\item[-] Use AI as a substitute for learning
\end{itemize}

\textbf{Remember:} The goal is to develop YOUR ability to solve these problems, with AI as a learning aid.
\end{tcolorbox}

\end{document}
